\noindent Мета роботи: дослідження сучасних програмних бібліотек реалізації криптографічних 
механізмів, критеріїв їх вибору для різних програмно-апаратних платформ та 
отримання практичних навичок коректного використання криптографічних API. 

\section{Дослідження принципів реалізації криптографічних механізмів}

\subsection{Чому при розробці прикладного програмного забезпечення не рекомендується самостійно реалізовувати стандартні криптографічні алгоритми?}

“Custom Algorithms → Don't do this”. \cite{OWASP}

Не рекомендується самостійно реалізовувати стандартні криптографічні алгоритми, оскільки навіть невелика помилка 
в їх реалізації може призвести до появи вразливості та зниження рівня захисту даних. Криптографічні бібліотеки вже містять 
готові реалізації поширених алгоритмів, які були протестовані та перевірені.

Крім того, відомі криптографічні бібліотеки регулярно оновлюються для виправлення виявлених помилок і вразливостей. 
Тому замість створення власної реалізації стандартного алгоритму доцільно використовувати надійні та широко перевірені 
криптографічні бібліотеки, як це рекомендує OWASP.

\subsection{Які основні групи криптографічних примітивів та допоміжних функцій зазвичай надають сучасні криптографічні бібліотеки?}

Сучасні криптографічні бібліотеки зазвичай містять такі основні групи криптографічних примітивів: \cite{Crypto Primitive Wiki}:

\begin{itemize}
\item односторонні хеш-функції --- обчислюють хеш-значення повідомлення;
\item симетрична криптографія --- забезпечує шифрування та розшифрування за допомогою одного ключа;
\item криптографія з відкритим ключем --- використовує різні ключі для шифрування та розшифрування;
\item цифрові підписи --- використовуються для підтвердження автентичності автора та цілісності повідомлення;
\item криптографічно стійкі генератори псевдовипадкових чисел (CSPRNG) --- використовуються для отримання непередбачуваних випадкових значень, 
зокрема криптографічних ключів та токенів.
\end{itemize}

Крім наведених примітивів, сучасні криптографічні бібліотеки можуть надавати й додаткові механізми, зокрема функції узгодження та отримання ключів, MAC, 
а також засоби роботи із сертифікатами та ключами.


\subsection{У чому полягає різниця між криптографічним примітивом та криптографічним протоколом? Навести приклади.}

Криптографічний примітив --- низькорівневий криптографічний алгоритм, який використовується як базовий будівельний
блок для криптографічних алгоритмів вищого рівня \cite{Cryptographic Primitive}.

Криптографічний протокол --- послідовність кроків та обмінів повідомленнями між кількома сторонами 
з метою досягнення певної цілі безпеки \cite{Crypto Protocols}.

Таким чином, основна відмінність полягає в тому, що криптографічний примітив є окремим алгоритмічним компонентом, тоді як криптографічний протокол описує 
взаємодію між сторонами та порядок використання криптографічних примітивів для забезпечення необхідних властивостей безпеки.

Базовими криптографічними примітивами є односторонні хеш-функції, наприклад SHA-256, симетричні криптографічні алгоритми, наприклад AES, 
асиметричні криптографічні алгоритми, наприклад RSA, а також алгоритми цифрового підпису. \cite{Crypto Primitive Wiki}
Прикладами криптографічних протоколів є Needham–Schroeder, Otway–Rees та Kerberos. \cite{Crypto Protocols}

\subsection{Що таке криптографічно стійкий генератор псевдовипадкових чисел (CSPRNG) та чому звичайні генератори псевдовипадкових чисел не повинні використовуватися для генерації криптографічних ключів?}

Випадкові числа та рядки використовуються для виконання важливих з погляду безпеки операцій, 
зокрема для генерації криптографічних ключів, векторів ініціалізації (ВІ), ідентифікаторів сеансів, 
CSRF-токенів та токенів для скидання пароля. Тому важливо, щоб такі значення було складно передбачити або відтворити зловмиснику.

Звичайні генератори псевдовипадкових чисел (ГПЧ) призначені переважно для завдань, де криптографічна безпека не потрібна. 
Вони можуть бути швидкими та зручними для таких операцій, як перемішування даних або випадковий вибір елементів. 
Проте їхній алгоритм або внутрішній стан у деяких випадках може бути передбачений, тому використання таких генераторів для створення 
криптографічних ключів може призвести до компрометації ключа.

Криптографічно стійкий генератор псевдовипадкових чисел (CSPRNG) призначений для генерування непередбачуваних випадкових значень, 
які можуть використовуватися в операціях, що потребують високого рівня безпеки. На відміну від звичайних ГПЧ, CSPRNG використовує механізми, 
які забезпечують криптографічну стійкість отриманих значень. Тому саме CSPRNG слід використовувати для генерації криптографічних ключів, токенів, 
ВІ та інших значень, що мають бути непередбачуваними. \cite{OWASP}


\subsection{У чому полягає різниця між низькорівневими (low-level) та високорівневими (high-level) криптографічними API? Які ризики може створювати неправильне використання низькорівневого API?}

Високорівневий API (High-level): створений за принципом <<захисту від помилок>> (\textit{secure-by-default}) \cite{Tink}. Він приховує складні налаштування й надає прості функції: розробник передає дані та ключ, а бібліотека сама генерує випадковий nonce/IV потрібної довжини, додає автентифікаційний тег і правильно обирає безпечний режим (наприклад, модуль \texttt{Fernet} у бібліотеці \texttt{cryptography}) \cite{cryptography}.
 
Низькорівневий API (Low-level): дає прямий доступ до базових примітивів (наприклад, шар \texttt{hazmat}) \cite{cryptography, cryptography git}. Тут розробник повинен вручну обирати розмір блоку, режим шифрування, спосіб доповнення (padding) і сам керувати векторами ініціалізації.


Ризики неправильного використання низькорівневого API \cite{OWASP, Libsodium}:
\begin{itemize}
    \item Повторне використання Nonce/IV: якщо в режимах на зразок AES-GCM або ChaCha20 зашифрувати два різні повідомлення з однаковим nonce на одному ключі, це призводить до витоку відкритого тексту та можливості підробки повідомлень;
    \item Вибір небезпечного режиму: розробник може помилково обрати застарілий режим на зразок ECB, який не приховує структуру даних;
    \item Відсутність перевірки цілісності: використання звичайного шифрування без автентифікації (без MAC або AEAD) дозволяє зловмиснику непомітно змінювати байти шифротексту;
    \item Атаки на доповнення (Padding Oracle): за неправильної ручної обробки помилок доповнення PKCS\#7 у режимі CBC зловмисник може повністю розшифрувати трафік.
\end{itemize}


\subsection{Що означає реалізація криптографічної операції з постійним часом виконання (constant-time implementation) та від яких атак вона повинна захищати?}

Реалізація з постійним часом означає, що криптографічна операція виконується за однакову кількість тактів процесора незалежно від того, які саме секретні дані (ключі, паролі, відкритий текст) обробляються \cite{Libsodium}.

Щоб цього досягти, у коді уникають:
\begin{itemize}
    \item умовних розгалужень (\texttt{if}, \texttt{switch}), які залежать від секретних бітів (бо процесор передбачає переходи з різною швидкістю);
    \item звернень до пам'яті за індексами, що залежать від секрету (щоб час операції не залежав від потрапляння в кеш-пам'ять процесора).
\end{itemize}

Захищає від:
\begin{itemize}
    \item Часові атаки (Timing attacks): коли зловмисник із високою точністю вимірює час відповіді системи. Наприклад, якщо порівнювати MAC звичайним циклом або оператором \texttt{==}, функція зупиняється на першому ж неправильному байті --- це дозволяє атакуючому байт за байтом підібрати правильне значення;
    \item Атаки через процесорний кеш (Cache-timing attacks): захищає від витоку інформації про те, які саме комірки пам'яті чи таблиці підстановок завантажувалися в кеш під час обчислень.
\end{itemize}


\subsection{Що таке FIPS 140-3 та в яких випадках наявність відповідної сертифікації криптографічного модуля може бути важливою?}

FIPS 140-3 є офіційним стандартом безпеки США та Канади (розроблений NIST), який встановлює вимоги до апаратних і програмних криптографічних модулів \cite{FIPS140-3, CMVP}. Він оцінює якість генерації ключів, захист від витоків, коректність реалізації алгоритмів та механізми самотестування модуля. Стандарт поділяється на 4 рівні безпеки (від базового програмного захисту Level 1 до серйозного фізичного захисту від розтину Level 4) \cite{SP800-140}.

Наявність сертифіката важлива, коли:
\begin{itemize}
    \item Державні та оборонні контракти: у США/НАТО використання модулів без сертифікації FIPS заборонено на законодавчому рівні (вимога FISMA та FedRAMP);
    \item Фінансовий сектор і банки: міжнародні стандарти захисту платіжних систем (наприклад, PCI-DSS) вимагають використання сертифікованих криптографічних засобів або HSM для обробки банківських карток;
    \item Хмарні сервіси та B2B-продукти: сертифікація підтверджує для корпоративних клієнтів, що бібліотека не містить грубих помилок і пройшла незалежне тестування в акредитованій лабораторії.
\end{itemize}


\subsection{Чому для реалізації RSA та деяких інших асиметричних криптосистем необхідна арифметика багаторозрядних чисел? Які вимоги, крім швидкодії, висуваються до її реалізації у криптографічному програмному забезпеченні?}

Стандартний рівень стійкості RSA вимагає використання ключів довжиною від 2048 до 4096 біт \cite{SP800-57}. Звичайні процесори вміють напряму працювати лише з числами розрядністю 32 або 64 біти. Тому для RSA потрібні спеціальні бібліотечні алгоритми <<великих чисел>> (\textit{BigNum}), які розбивають гігантські числа на масиви та виконують над ними модульне множення й піднесення до степеня \cite{RFC8017}:
\begin{equation}
c \equiv m^e \pmod{N}, \quad m \equiv c^d \pmod{N}
\end{equation}

Вимоги до реалізації (крім швидкодії):
\begin{itemize}
    \item Захист від витоків за часом: модульне піднесення до степеня закритим ключем має виконуватися за однаковий час незалежно від кількості нулів чи одиниць у ключі (захист від часових атак);
    \item Криптографічне засліплення (Blinding): перемішування вхідного повідомлення з випадковим числом перед обчисленням закритим ключем ($m' = m \cdot r^e \pmod{N}$), щоб сторонній спостерігач не міг зіставити вхідні дані з часом або споживанням енергії \cite{RFC8017};
    \item Безпечне очищення пам'яті (Zeroization): проміжні результати та фрагменти ключа в оперативній пам'яті мають надійно затиратися нулями після завершення операції, щоб компілятор не оптимізував це очищення (\texttt{explicit\_bzero});
    \item Контроль переповнення (Memory Safety): гарантована безпека коду від помилок роботи з пам'яттю (buffer overflow) під час розрахунку великих добутків.
\end{itemize}
